% =====================================================
% 题型：解三角形角平分线最值解答题 (q_lyb_01_16_triangle_bisector.tex)
% =====================================================
\newcommand{\qLYBZeroOneSixteenStem}{%
  设 \(\triangle ABC\) 的内角 \(A, B, C\) 的对边分别为 \(a, b, c\)，\(\angle B\) 的角平分线 \(BD\) 与 \(AC\) 交于点 \(D\)。
  \begin{enumerate}[label=(\arabic*) , leftmargin=2em, itemsep=0.1em]
    \item 若 \(\vec{AD} = \dfrac{1}{5}\vec{AC}\)，求 \(\sin A : \sin C\) 的值；
    \item 若 \(b\cos C + \sqrt{3}b\sin C = a + c\)，且 \(BD = 2\)，求 \(c + 9a\) 的最小值。
  \end{enumerate}%
}

\newcommand{\qLYBZeroOneSixteenFigure}[1][1.3]{%
  \begin{tikzpicture}[scale=#1, >=stealth]
    % Coordinates of triangle
    \coordinate (B) at (0, 1.2);
    \coordinate (A) at (-1.0, 0);
    \coordinate (C) at (2.0, 0);
    
    % BD is the angle bisector.
    % D divides AC. Since AD = 1/5 AC => D is at 1/5 from A to C
    \coordinate (D) at ($(A)!0.2!(C)$);
    
    \draw[thick] (A) -- (B) -- (C) -- cycle;
    \draw[thick] (B) -- (D);
    
    \node[below left] at (A) {\small $A$};
    \node[above] at (B) {\small $B$};
    \node[below right] at (C) {\small $C$};
    \node[below] at (D) {\small $D$};
  \end{tikzpicture}%
}

\newcommand{\qLYBZeroOneSixteenAnswer}{%
  (1) \(\sin A : \sin C = 4:1\)；\\
  (2) \(c + 9a\) 的最小值为 \(\dfrac{32\sqrt{3}}{3}\)。%
}

\newcommand{\qLYBZeroOneSixteenSolution}{%
  (1) 因为 \(\vec{AD} = \dfrac{1}{5}\vec{AC}\)，所以 \(AD = \dfrac{1}{5}AC\)，从而 \(AD : DC = 1 : 4\)。
  因为 \(BD\) 是 \(\angle B\) 的角平分线，由三角形角平分线定理得：
  \(\dfrac{AB}{BC} = \dfrac{AD}{DC} = \dfrac{1}{4} \implies \dfrac{c}{a} = \dfrac{1}{4}\)。
  根据正弦定理，有：
  \(\dfrac{\sin A}{\sin C} = \dfrac{a}{c} = 4\)。
  所以 \(\sin A : \sin C\) 的值为 \(4:1\)。\\
  
  (2) 在 \(\triangle ABC\) 中，由正弦定理将已知等式 \(b\cos C + \sqrt{3}b\sin C = a + c\) 边化为角，得：
  \(\sin B \cos C + \sqrt{3}\sin B \sin C = \sin A + \sin C\)。
  因为 \(A = \pi - (B+C)\)，所以 \(\sin A = \sin(B+C) = \sin B \cos C + \cos B \sin C\)。
  代入上式化简，得：
  \(\sin B \cos C + \sqrt{3}\sin B \sin C = \sin B \cos C + \cos B \sin C + \sin C \implies \sqrt{3}\sin B \sin C = \cos B \sin C + \sin C\)。
  因为 \(C \in (0, \pi) \implies \sin C > 0\)，等式两边同除以 \(\sin C\)，得：
  \(\sqrt{3}\sin B - \cos B = 1 \implies 2\sin\left(B - \dfrac{\pi}{6}\right) = 1 \implies \sin\left(B - \dfrac{\pi}{6}\right) = \dfrac{1}{2}\)。
  因为 \(B \in (0, \pi) \implies B - \dfrac{\pi}{6} \in \left(-\dfrac{\pi}{6}, \dfrac{5\pi}{6}\right)\)，
  所以 \(B - \dfrac{\pi}{6} = \dfrac{\pi}{6} \implies B = \dfrac{\pi}{3}\)。\\
  因为 \(BD = 2\) 是 \(\angle B\) 的平分线，根据面积关系 \(S_{\triangle ABC} = S_{\triangle ABD} + S_{\triangle CBD}\)，有：
  \(\dfrac{1}{2}ac \sin B = \dfrac{1}{2}c \cdot BD \sin\dfrac{B}{2} + \dfrac{1}{2}a \cdot BD \sin\dfrac{B}{2}\)。
  代入 \(B = \dfrac{\pi}{3}\)，\(\sin B = \dfrac{\sqrt{3}}{2}\)，\(\sin\dfrac{B}{2} = \sin\dfrac{\pi}{6} = \dfrac{1}{2}\)，以及 \(BD = 2\)，得：
  \(\dfrac{\sqrt{3}}{4}ac = \dfrac{1}{2}c + \dfrac{1}{2}a \implies \dfrac{\sqrt{3}}{2}ac = a + c\)。
  两边同除以 \(ac\)，可得 \(\dfrac{1}{a} + \dfrac{1}{c} = \dfrac{\sqrt{3}}{2}\)。\\
  要求 \(c + 9a\) 的最小值，利用常数代换法：
  \(c + 9a = (c + 9a) \times \dfrac{2}{\sqrt{3}}\left(\dfrac{1}{a} + \dfrac{1}{c}\right) = \dfrac{2}{\sqrt{3}} \left( \dfrac{c}{a} + 1 + 9 + \dfrac{9a}{c} \right) = \dfrac{2}{\sqrt{3}}\left(10 + \dfrac{c}{a} + \dfrac{9a}{c}\right)\)。
  由均值不等式 \(\dfrac{c}{a} + \dfrac{9a}{c} \ge 2\sqrt{\dfrac{c}{a} \times \dfrac{9a}{c}} = 6\)，
  所以 \(c + 9a \ge \dfrac{2}{\sqrt{3}}(10 + 6) = \dfrac{32}{\sqrt{3}} = \dfrac{32\sqrt{3}}{3}\)。
  当且仅当 \(\dfrac{c}{a} = \dfrac{9a}{c} \implies c = 3a\) 时等号成立。
  此时由 \(\dfrac{1}{a} + \dfrac{1}{3a} = \dfrac{\sqrt{3}}{2} \implies \dfrac{4}{3a} = \dfrac{\sqrt{3}}{2} \implies a = \dfrac{8\sqrt{3}}{9}\)，符合条件。
  所以 \(c + 9a\) 的最小值为 \(\dfrac{32\sqrt{3}}{3}\)。%
}
